\chapter{Introduzione}
\label{cap:introduzione}

\section{Contesto e motivazione}

La vista consente di acquisire rapidamente una grande quantità di informazioni
sull'ambiente circostante: permette di riconoscere oggetti e persone,
interpretare segnali ed etichette, leggere testi e comprendere la disposizione
degli elementi nello spazio. Quando questa capacità è ridotta o assente,
l'accesso a tali informazioni può richiedere strumenti assistivi oppure il
supporto di altre persone.

La disabilità visiva comprende condizioni differenti e non corrisponde a
un'unica esperienza percettiva. Può interessare la visione da vicino o da
lontano e presentarsi con livelli diversi di gravità. Alcune persone conservano
una capacità visiva parziale, mentre altre non possono utilizzare la vista per
accedere alle informazioni presenti nell'ambiente. Con il termine
\emph{ipovisione} si indica, in questo lavoro, una condizione nella quale è
presente un residuo visivo, ma la percezione disponibile può non essere
sufficiente per riconoscere con affidabilità dettagli, testi o oggetti in ogni
condizione. Anche fattori come illuminazione, contrasto e organizzazione dello
spazio possono influire sull'autonomia nelle attività quotidiane. Secondo
l'Organizzazione Mondiale della Sanità, almeno 2,2 miliardi di persone nel
mondo presentano una compromissione della visione da vicino o da lontano
\cite{who_blindness_vision_impairment_2026}.

In questo contesto, le tecnologie assistive possono contribuire a rendere
accessibili attraverso il canale uditivo informazioni originariamente
disponibili soltanto in forma visiva. I progressi nella computer vision, nei
modelli visivo-linguistici e nella sintesi vocale rendono oggi possibile
analizzare un'immagine, produrne una descrizione testuale e convertirla in
parlato. Affinché questo processo sia realmente utile, tuttavia, non è
sufficiente generare una descrizione: occorre stabilire quando fornirla, quale
elemento descrivere e come evitare informazioni ripetitive o poco pertinenti.

% Qui puoi aggiungere, se lo desideri, una breve motivazione personale
% che spieghi perché hai scelto questo specifico problema di tesi.

\section{Problema di ricerca}

L'ambiente quotidiano è dinamico e contiene una quantità elevata di
informazioni visive. Una descrizione continua e indiscriminata rischierebbe di
sovraccaricare l'utente, sovrapporre più messaggi vocali e ripetere contenuti
già comunicati. Al contrario, un sistema attivabile soltanto mediante comandi
espliciti richiederebbe all'utente di formulare ogni volta una richiesta,
anche quando avrebbe bisogno semplicemente di conoscere un cambiamento
avvenuto nelle vicinanze.

Si presenta quindi la necessità di conciliare due esigenze complementari:
mantenere una consapevolezza generale di ciò che accade nell'ambiente e
consentire all'utente di richiedere informazioni mirate su un elemento
specifico. Il sistema deve inoltre limitare la latenza, evitare descrizioni
ridondanti e comunicare soltanto informazioni che possano essere formulate con
un livello adeguato di affidabilità.

La domanda che guida il presente lavoro può essere formulata nel modo
seguente:

\begin{quote}
\emph{È possibile realizzare un sistema browser-based che fornisca
descrizioni vocali contestuali, sia automaticamente sia attraverso un gesto di
puntamento, senza richiedere hardware proprietario?}
\end{quote}

\section{Obiettivi del progetto}

L'obiettivo della tesi è la progettazione e lo sviluppo di VisionCaption, un
sistema di visione assistita rivolto a persone cieche e ipovedenti, capace di
generare descrizioni vocali della scena osservata dalla fotocamera di un
dispositivo. L'applicazione è concepita per funzionare tramite browser, così da
poter essere utilizzata su dispositivi comuni, come computer, smartphone e
tablet, senza dipendere da uno specifico visore o da sensori proprietari.

Le descrizioni devono essere brevi, comprensibili e formulate in lingua
italiana. L'obiettivo non è sostituire gli ausili tradizionali per
l'orientamento e la mobilità, né realizzare un dispositivo medico, ma fornire
un ulteriore canale di accesso alle informazioni visive. Il sistema è quindi
pensato come supporto alla comprensione del contesto e non come fonte unica
sulla quale basare decisioni relative alla sicurezza personale.

Per rispondere a esigenze differenti, VisionCaption prevede due modalità di
interazione: una modalità automatica, orientata alla consapevolezza generale
dell'ambiente, e una modalità di puntamento, destinata all'esplorazione
intenzionale di un oggetto.

\section{Modalità di interazione}

\subsection{Modalità automatica}

La modalità automatica non richiede un intervento esplicito da parte
dell'utente. Il sistema osserva il flusso proveniente dalla fotocamera e produce
una nuova descrizione quando rileva un cambiamento sufficientemente rilevante
nella scena. In questo modo l'utente può ricevere informazioni su ciò che lo
circonda durante le normali attività quotidiane, evitando che una scena
sostanzialmente invariata generi continuamente la stessa caption.

Questa modalità privilegia quindi la continuità dell'assistenza, ma deve
mantenere un equilibrio tra tempestività e quantità delle informazioni. Una
frequenza eccessiva renderebbe infatti l'audio invasivo, mentre una frequenza
troppo bassa potrebbe omettere cambiamenti utili.

\subsection{Modalità di puntamento}

La modalità di puntamento offre un'interazione più intenzionale. L'utente può
indicare con il dito l'elemento sul quale desidera ricevere informazioni,
senza dover formulare verbalmente una domanda o selezionare manualmente una
regione dello schermo. Il sistema riconosce una posa di puntamento mantenuta in
modo stabile, interpreta la direzione indicata e genera una descrizione
specifica del possibile referente.

Questa modalità può essere impiegata, ad esempio, per conoscere la natura di
un oggetto, distinguerlo dagli elementi vicini oppure leggere il testo visibile
su un'etichetta o una confezione. Il gesto deittico costituisce così un
collegamento diretto tra l'intenzione dell'utente e la porzione di ambiente da
descrivere.

\section{Contributo del lavoro}

Il contributo del progetto non consiste nella definizione di un nuovo modello
fondazionale di intelligenza artificiale, ma nell'integrazione di tecniche di
visione artificiale, modelli multimodali e sintesi vocale all'interno di
un'unica esperienza assistiva. L'aspetto centrale è la coesistenza di due
strategie complementari: una descrizione automatica sensibile ai cambiamenti
della scena e una descrizione attivata da un gesto naturale di puntamento.

La scelta di un'interfaccia browser-based mira inoltre a ridurre la dipendenza
da piattaforme e dispositivi specifici. Il lavoro affronta pertanto non solo il
riconoscimento del contenuto visivo, ma anche problemi legati alla selezione
dell'informazione, alla latenza, alla ripetizione delle caption e alla
semplicità dell'interazione.

\section{Organizzazione della tesi}

Il presente capitolo introduce il contesto, il problema e gli obiettivi del
lavoro. Il Capitolo~\ref{cap:tecnologie} presenta le principali tecnologie
utilizzate, mentre il capitolo dedicato allo stato dell'arte confronta
VisionCaption con le soluzioni assistive esistenti. I capitoli successivi
descrivono la progettazione del sistema, le scelte implementative e la
valutazione dei risultati ottenuti.

% Aggiorna l'ultimo periodo quando l'indice definitivo dei capitoli sarà
% stabilito, indicando esplicitamente il contenuto di ciascun capitolo.
